\documentclass[10pt]{scrartcl}
\usepackage[a4paper, total={5.5in, 8in}]{geometry}
\usepackage[usenames,dvipsnames,svgnames]{xcolor}
\usepackage[shortlabels]{enumitem}
\usepackage[framemethod=TikZ]{mdframed}
\usepackage{amsmath,amssymb,amsthm}
\usepackage[colorlinks]{hyperref}

\hypersetup{
	colorlinks=true,
	linkcolor=blue,
	filecolor=blue,      
	urlcolor=blue,
	pdftitle={MAT 487 Spr25}
}

\usepackage{microtype}
\usepackage{mathtools}
\usepackage[headsepline]{scrlayer-scrpage}
\usepackage{thmtools}
\usepackage[textsize=small]{todonotes}

\addtolength{\textheight}{3.14cm}
\providecommand{\ol}{\overline}
\providecommand{\eps}{\varepsilon}
\providecommand{\half}{\frac{1}{2}}
\providecommand{\dang}{\measuredangle} %% Directed angle
\providecommand{\CC}{\mathbb C}
\providecommand{\EE}{\mathbb E}
\providecommand{\FF}{\mathbb F}
\providecommand{\NN}{\mathbb N}
\providecommand{\QQ}{\mathbb Q}
\providecommand{\RR}{\mathbb R}
\providecommand{\ZZ}{\mathbb Z}
\providecommand{\dg}{^\circ}
\providecommand{\ind}[1]{\mathbf 1_{#1}}
\providecommand{\ii}{\item}
\providecommand{\alert}[1]{{\sffamily\textbf{\textcolor{magenta}{#1}}}}
\providecommand{\opname}{\operatorname}
\providecommand{\li}{\opname{li}}
\providecommand{\ls}{\opname{ls}}
\providecommand{\inv}{^{-1}}
\providecommand{\setdiff}{\smallsetminus}
\providecommand{\mb}[1]{\mathbf{#1}}
\providecommand{\abs}[1]{\left|{#1}\right|}
\providecommand{\symdiff}{\Delta}
\providecommand{\jim}{m_{*,(J)}}
\providecommand{\jom}{m^{*,(J)}}
\providecommand{\pcint}{\mathrm{p.c.}\int}
\providecommand{\dif}{\;d}
\providecommand{\dist}{\operatorname{dist}}
\providecommand{\diam}{\operatorname{diam}}
\renewcommand{\hat}{\widehat}
\providecommand{\norm}[1]{\| #1 \|}

\reversemarginpar
\setlength{\marginparwidth}{3cm}

\mdfdefinestyle{mdbluebox}{roundcorner=10pt,innerbottommargin=9pt,
	linecolor=blue,backgroundcolor=TealBlue!5,}
\declaretheoremstyle[headfont=\sffamily\bfseries\color{MidnightBlue},
mdframed={style=mdbluebox},]{thmbluebox}

\mdfdefinestyle{mdredbox}{frametitlefont=\bfseries,innerbottommargin=8pt,
	nobreak=true,backgroundcolor=Salmon!5,linecolor=RawSienna,}
\declaretheoremstyle[headfont=\bfseries\color{RawSienna},
mdframed={style=mdredbox},headpunct={\\[3pt]},postheadspace=0pt,]{thmredbox}

\mdfdefinestyle{mdgreenbox}{linecolor=ForestGreen,backgroundcolor=ForestGreen!5,
	linewidth=2pt,rightline=false,leftline=true,topline=false,bottomline=false,}
\declaretheoremstyle[headfont=\bfseries\sffamily\color{ForestGreen!70!black},
mdframed={style=mdgreenbox},headpunct={ --- },]{thmgreenbox}

\mdfdefinestyle{mdorangebox}{linecolor=RedOrange,backgroundcolor=RedOrange!5,
	linewidth=2pt,rightline=false,leftline=true,topline=false,bottomline=false,}
\declaretheoremstyle[headfont=\bfseries\sffamily\color{RedOrange!70!black},
mdframed={style=mdorangebox},headpunct={ --- },]{thmorangebox}

\mdfdefinestyle{mdblackbox}{linecolor=black,backgroundcolor=RedViolet!5!gray!5,
	linewidth=3pt,rightline=false,leftline=true,topline=false,bottomline=false,}
\declaretheoremstyle[mdframed={style=mdblackbox}]{thmblackbox}


\declaretheorem[style=thmredbox,name=Problem,numberwithin=section]{problem}

\declaretheorem[style=thmredbox,name=Example,sibling=problem]{example}
\declaretheorem[style=thmredbox,name=$\bigstar$ Problem,sibling=problem]{niceprob}

\declaretheorem[style=thmbluebox,name=Theorem,sibling=problem]{theorem}
\declaretheorem[style=thmbluebox,name=Corollary,sibling=theorem]{corollary}
\declaretheorem[style=thmbluebox,name=Proposition,sibling=theorem]{prop}
\declaretheorem[style=thmbluebox,name=Lemma,sibling=theorem]{lemma}
\declaretheorem[style=thmbluebox,name=Theorem,numbered=no]{theorem*}
\declaretheorem[style=thmbluebox,name=Lemma,numbered=no]{lemma*}
\declaretheorem[style=thmgreenbox,name=Claim,numberwithin=theorem]{claim}
\declaretheorem[style=thmgreenbox,name=Claim,numbered=no]{claim*}
\declaretheorem[style=thmblackbox,name=Remark,sibling=theorem]{remark}
\declaretheorem[style=thmblackbox,name=Remark,numbered=no]{remark*}
\declaretheorem[style=thmgreenbox,name=Definition,sibling=theorem]{definition}
\declaretheorem[style=thmgreenbox,name=Definition,numbered=no]{definition*}
\declaretheorem[style=thmorangebox,name=Warning,sibling=theorem]{warning}
\declaretheorem[style=thmorangebox,name=Warning,numbered=no]{warning*}
\declaretheorem[style=thmorangebox,name=Note,numbered=no]{note}
\declaretheorem[style=thmblackbox,name=Exercise,sibling=problem]{exercise}
\declaretheorem[style=thmblackbox,name=Exercise,numbered=no]{exercise*}

\newenvironment{soln}{\begin{proof}[Solution]}{\end{proof}}
\newenvironment{walkthrough}{\noindent \textbf{Walkthrough.}}{\medskip}
\newlist{walk}{enumerate}{3}
\setlist[walk]{label=\bfseries (\alph*)}

\providecommand{\pow}{\operatorname{Pow}}
\providecommand{\vocab}[1]{{\textbf{\textcolor{ForestGreen}{#1}}}}
\providecommand{\lcm}{\operatorname{lcm}}
\providecommand{\ord}{\operatorname{ord}}
\providecommand{\ex}{\opname{ex}}
\providecommand{\floor}[1]{\left\lfloor #1\right\rfloor}
\providecommand{\ceil}[1]{\left\lceil #1\right\rceil}

\usepackage{asymptote}
\begin{asydef}
	defaultpen(fontsize(10pt));
	unitsize(1cm);
	size(8cm); // set a reasonable default
	usepackage("amsmath");
	usepackage("amssymb");
	settings.tex="pdflatex";
	settings.outformat="pdf";
	// Replacement for olympiad+cse5 which is not standard
	import geometry;
	// recalibrate fill and filldraw for conics
	void filldraw(picture pic = currentpicture, conic g, pen fillpen=defaultpen, pen drawpen=defaultpen)
	{ filldraw(pic, (path) g, fillpen, drawpen); }
	void fill(picture pic = currentpicture, conic g, pen p=defaultpen)
	{ filldraw(pic, (path) g, p); }
	// some geometry
	pair foot(pair P, pair A, pair B) { return foot(triangle(A,B,P).VC); }
	pair orthocenter(pair A, pair B, pair C) { return orthocentercenter(A,B,C); }
	pair centroid(pair A, pair B, pair C) { return (A+B+C)/3; }
	// cse5 abbreviations
	path CP(pair P, pair A) { return circle(P, abs(A-P)); }
	path CR(pair P, real r) { return circle(P, r); }
	pair IP(path p, path q) { return intersectionpoints(p,q)[0]; }
	pair OP(path p, path q) { return intersectionpoints(p,q)[1]; }
	path Line(pair A, pair B, real a=0.6, real b=a) { return (a*(A-B)+A)--(b*(B-A)+B); }
	// cse5 more useful functions
	picture CC() {
		picture p=rotate(0)*currentpicture;
		currentpicture.erase();
		return p;
	}
	pair MP(Label s, pair A, pair B = plain.S, pen p = defaultpen) {
		Label L = s;
		L.s = "$"+s.s+"$";
		label(L, A, B, p);
		return A;
	}
	pair Drawing(Label s = "", pair A, pair B = plain.S, pen p = defaultpen) {
		dot(MP(s, A, B, p), p);
		return A;
	}
	path Drawing(path g, pen p = defaultpen, arrowbar ar = None) {
		draw(g, p, ar);
		return g;
	}
\end{asydef}

\ihead{\footnotesize\textbf{MAT 487 Spring 2025}}
\ohead{\footnotesize \today}

\title{\vspace{-2.0cm} Roth's theorem, done thrice}
\author{\large Aditya Pahuja}
\date{\large Presented on 2025-02-26}
\begin{document}
\maketitle
\tableofcontents
\pagebreak

\section{Act 1: Finite field vector spaces}
\subsection{Fourier crash course}
We will start by proving a version of Roth's theorem in
$\FF_p^n$ for odd primes $p$, which will serve as a nice environment to develop
the technique of using Fourier analysis to prove theorems in
additive combinatorics.

If you are entirely unfamiliar with what the Fourier transform does go watch
\href{https://www.youtube.com/watch?v=spUNpyF58BY}{this 3Blue1Brown video}.
It will make your life better. I will do a poor attempt at explaining, though.

The goal of the Fourier transform is to take a function
and write it as a linear combination of exponential functions;
specifically, we want to know, ``what were the
frequencies $\omega$ of the exponentials $e^{i\omega t}$
that were linearly combined to make this function?''
The way it does this is explained very nicely in the video,
but the gist is that if you have your function $f$,
then its Fourier transform $\hat f(r)$, when given a candidate frequency $r$,
will generally be pretty small, but it will spike when the frequency $r$ is
``resonant'' with the function $f$: thus, we are often interested in isolating
the ``Fourier coefficients'' $\hat f(r)$ with large absolute value,
since we want to find the $r$ values that resonate with our function.

In this spirit, we will begin with our first definition:
\begin{definition}[Fourier transform, $\FF_p^n$]
	Let $f\colon\FF_p^n\to\CC$ be a function. The \vocab{Fourier transform}
	$\hat f\colon\FF_p^n\to\CC$ is the average
	\[\hat f(r) := \EE_{x\in\FF_p^n} f(x)\omega^{-r\cdot x}
	= \frac{1}{p^n}\sum_{x\in\FF_p^n} f(x)\omega^{-r\cdot x}\]
	where $r\cdot x = r_1x_1 + \cdots + r_nx_n$.
\end{definition}

Notably, $\hat f(0)$ is the average of $f$.

One should think about your standard functions $f$ as being
in ``physical space,'' say representing sound waves or some similar information,
while their transforms $\hat f$ are in ``frequency space,''
keeping track of how resonant $f$ is with each frequency.
Here ``resonance'' is something of a synonym for ``orthogonality.''

In accordance with our analogy of frequencies, we have the so-called
Fourier inversion formula:
\begin{theorem}[Fourier inversion]
	Let $f\colon\FF_p^n\to\CC$. We have the identity
	\[f(x) = \sum_{r\in\FF_p^n} \hat f(r)\omega^{r\cdot x}.\]
\end{theorem}

%\begin{remark}
%	The dichotomy between averaging in physical space
%	and summing in frequency space is a result of
%	the fact, given a compact abelian group $G$,
%	its dual group $\hat G$ consisting of homomorphisms $G\to\RR/\ZZ$
%	has the discrete topology when endowed with the so-called
%	\vocab{compact-open topology}, so $G$ is given a
%	``probability Haar measure'' (a translation-invariant probability
%	measure on $G$ that obeys outer and inner regularity,
%	i.e. is well-behaved with the topology of $G$)
%	while $\hat G$ is just given the counting measure.
%\end{remark}

Additionally, we have the following identity,
which we will shortly see exhibits the unitarity of the Fourier transform:
\begin{theorem}[Parseval/Plancherel]
	Given $f,g\colon\FF_p^n\to\CC$,
	\[\EE[\ol{ f(x)}g(x)] = \sum \ol{\hat f(r)} \hat g(r).\]
	In particular, when $f=g$,
	\[\EE[|f(x)|^2] = \sum|\hat f(r)|^2.\]
	Here, and going forward, the abbreviated $\EE h(x)$ and $\sum \hat h(r)$
	are both assumed to be uniformly over all $x,h\in\FF_p^n$.
\end{theorem}
For consistency we will refer to this theorem as just Parseval's theorem.

The point of all of this is that in physical space,
we have the averaging inner product and norm
\[\langle f,g\rangle := \EE[\ol{f(x)}g(x)],\quad
\norm{f}_2 := \langle f,f\rangle^{1/2},\]
while in frequency space, we have the summing inner product and norm
\[\langle \alpha,\beta\rangle_{\ell^2} := \sum\ol{\alpha(x)}\beta(x),\quad
\norm{\alpha}_{\ell^2} = \langle\alpha,\alpha\rangle_{\ell^2}^{1/2}.\]

With respect to either inner product, we then have the orthonormal basis
of characters $\gamma_r(x) := \omega^{r\cdot x}$ (check this!).
Thus the Fourier transform can be written as
\[\hat f(r) = \langle \gamma_r, f\rangle\]
and Parseval's identity directly states
\[\langle f,g\rangle = \langle \hat f,\hat g\rangle_{\ell^2},\quad
\norm f_2 = \norm{\hat f}_{\ell^2}.\]

With this stronger understanding of the structure of the Fourier transform,
proving the inversion formula and the Parseval identity is very simple:
\begin{proof}[Proof of Fourier inversion forumla]
	Observe that
	\[\sum_r \hat f(r)\omega^{r\cdot x} =
	\sum_r \langle \gamma_r,f\rangle\gamma_r = f.\qedhere\]
\end{proof}

\begin{proof}[Proof of Parseval]
	Splitting $f$ and $g$ into the coordinate sums and using
	the sesquilinearirty of the inner product, we can write
	\[\langle f,g\rangle = \sum_{r\in\FF_p^n}
	\langle f,\gamma_r\rangle\langle\gamma_r,g\rangle
	= \sum_r \ol{\hat f(r)}\hat g(r) = \langle \hat f,\hat g\rangle_{\ell^2}.
	\qedhere\]
\end{proof}

Another important operation to keep in mind is the convolution:
\begin{definition}[Convolution]
	The \vocab{convolution} of two functions $f,g\colon\FF_p^n\to\CC$ is
	the function $f*g\colon\FF_p^n\to\CC$ defined by
	\[(f*g)(x) = \EE_{a+b=x}[f(a)g(b)].\]
\end{definition}

\begin{example}
	The convolution can be thought of as a ``smoothing'' operation:
	for example, if $W$ is a subspace of $\FF_p^n$ and $\mu_W = (p^n/|W|)\ind W$
	the normalized indicator function of $W$ (so that its average is $1$),
	then $f*\mu_W(x)$ is obtained by replacing the average of $f$ at $x$
	with its average on the coset $x+W$.
\end{example}

\begin{exercise}
	Let $A$ and $B$ be two subsets of $\FF_p^n$. What can you say about
	$\ind A*\ind B$?
\end{exercise}

The reason convolution is so important in this context
is because of how incredibly nicely it behaves with the Fourier transform:
\begin{theorem}[Convolutions become products]
	For any $f,g\in\FF_p^n\to\CC$ and $r\in\FF_p^n$,
	\[\hat{f*g}(r) = \hat f(r)\hat g(r).\]
\end{theorem}

\begin{proof}
	Omitted; it's just some algebraic manipulation.
\end{proof}

Of course, we can generalize this to $(f_1*\cdots*f_k)^\wedge =
\hat f_1\hat f_2\cdots\hat f_k$ (where we write $f^\wedge$ instead of
$\hat f$ to avoid hideous typesetting).

\subsection{Roth's theorem}
In the context of Roth's theorem --- or rather, 3-APs in general ---
we often want to think about the quantity
\[\Lambda(f,g,h) := \EE_{x,y}f(x)g(x+y)h(x+2y).\]
The special case $f=g=h$ appears often enough as well
that we give it a separate name:
\[\Lambda_3(f) := \Lambda(f,f,f).\]
In particular
\[\Lambda_3(\ind A) = p^{-2n}\abs{\{(x,y):x,x+y,x+2y\in A\}}\]
is the proportion of pairs $(x,y)$ that induce an arithmetic progression in $A$,
which can be thought of as the 3-AP density of $A$;
note that we include trivial 3-APs with $y=0$.

The punchline of all of this is that we can write the condition
``$x$, $y$, $z$ form an arithmetic sequence'' as precisely $x-2y+z=0$,
which means that we get
\begin{align*}
	\Lambda(f,g,h) &=\EE_{x,y}f(x)g(x+y)h(x+2y)\\
	&= \EE_{x,y,z:\;x-2y+z=0} f(x)g(y)h(z)\\
	&= \EE_{x,y',z:\; x+y+z=0} f(x)g_1(y)h(z)\\
	&= (f*g_1*h)(0)\\
	&= \sum_r \hat{f*g_1*h}(r)\omega^{0\cdot r}\\
	&= \sum_r \hat f(r)\hat{g_1}(r)\hat h(r)\\
	&= \sum_r \hat f(r)\hat g(-2r)\hat h(r)
\end{align*}
where $g_1(y) = g(-y/2)$; note that here we used the fact that
$y\mapsto -2y$ is a bijection on $\FF_p^n$.
To reiterate, $\Lambda(f,g,h)$, an averaging quantity,
can be converted, in the Fourier spirit, to a sum of products of
$\hat f$, $\hat g$, $\hat h$ (and, as a bonus, reduces the number of
parameters from $2$ to $1$ in the process).
Hopefully this is all convincing you that $\Lambda$ is really quite
a useful operation to think about.

Let us now state Roth's theorem in the finite field setting:
\begin{theorem}[Roth's theorem, mk. I]
	Every 3-AP-free subset of $\FF_p^n$ has size $O(p^n/n)$.
\end{theorem}

The strategy of each proof we will give today is very similar to the
graph-theoretic edition of this proof
(but we have our own fun details, of course):
\begin{enumerate}[(1)]
	\ii A 3-AP-free set has a large Fourier coefficient.
	\ii A large Fourier coefficient implies density increment
	on some hyperplane.
	\ii Restrict to the hyperplane and iterate.
\end{enumerate}

Understanding the nature of how quickly this algorithm halts
will ultimately give us a bound on the size of the 3-AP-free set.
(When we say that a set $A$ has a large Fourier coefficient,
we are committing an abuse of language and identifying $A$ with $\ind A$.)
As we mentioned earlier, having a large Fourier coefficient
$\hat{\ind{A}}(r)$ for $r\ne 0$ implies that there is some nonuniformity
in the direction of $r$, while having a small coefficient implies that
$A$ looks roughly random in the direction of $r$.

In accordance with our strategy, let us begin by finding a 3-AP-free set
with large Fourier coefficient. The starting point
is a counting lemma, which gives us a bound involving 3-AP density
that we can exploit.
\begin{lemma}[Counting lemma]
	Let $f\colon\FF_p^n\to[0,1]$. Then,
	\[\abs{\Lambda_3(f) - (\EE f)^3} \le \max_{r\ne 0}|\hat f(r)|\norm f_2^2.\]
\end{lemma}
Intuitively, this is saying that $\Lambda_3(f)$ should be
pretty close to $\EE f$ --- that is, $x$, $x+y$, and $x+2y$
behave roughly independently --- unless there is a particularly
resonant vector $r\in\FF_p^n$.

\begin{proof}[Proof of counting lemma]
	From the identity on $\Lambda(f,g,h)$, we have
	\[\Lambda_3(f) = \sum_r \hat f(r)^3
	= \hat f(0)^3 + \sum_{r\ne 0} \hat f(r)^2\hat f(-2r)
	= (\EE f)^3 + \sum_{r\ne 0} \hat f(r)^2\hat f(-2r).\]
	Taking absolute values,
	we can then get rid of the $f(-2r)$ by bounding the expression with
	\[|\Lambda_3(f) - (\EE f)^3| \le \sum_{r\ne 0} |\hat f(r)|^2|\hat f(-2r)|
	\le \max_{r\ne 0}|\hat f(r)|\cdot \sum_{r}\hat f(r)^2
	= \max_{r\ne 0}|\hat f(r)|\cdot \norm f_2^2.\qedhere\]
\end{proof}

The large Fourier coefficient now follows:
\begin{lemma}[Large Fourier coefficient]
	Let $A\subseteq\FF_p^n$ be 3-AP-free and $\alpha = |A|/p^n$.
	If $p^n \ge 2\alpha^{-2}$, then there is $r\ne0$ such that
	$|\hat {\ind A}(r)|\ge\frac{\alpha^2}2$.
\end{lemma}

\begin{proof}
	We start by noting that the only 3-APs counted by $\Lambda_3(\ind A)$
	are the trivial ones, of which there are $|A|$; thus this quantity is
	$|A|/p^{2n} = \alpha/p^n$.
	Applying the counting lemma to $\ind A$, we have
	\[\frac{\alpha^3}{2}\le
	\alpha^3 - \frac\alpha{p^n} =\alpha^3 - \Lambda_3(\ind A)
	\le \max_{r\ne 0}|\hat{\ind A}(r)|\cdot \norm{\ind A}_2^2
	= \max_{r\ne 0}|\hat{\ind A}(r)|\cdot\alpha,\]
	so dividing by $\alpha$ gives us what we want.
\end{proof}

To reiterate, we have now shown that, as long as a 3-AP-free set
has small enough density, we can find a large Fourier coefficient.
Now, we will restrict to a hyperplane, noting that hyperplanes
are just copies of $\FF_p^{n-1}$ embedded as subspace of $\FF_p^n$.
Specifically, we will restrict to a hyperplane
whose normal vector is most resonant with our set $A$, namely
one orthogonal to the $r$ we found via the previous lemma.

\begin{lemma}[Density increment]
	Let $A\subseteq \FF_p^n$ with $\alpha = |A|/p^n$.
	Suppose that $|\hat{\ind A}(r)| \ge\delta > 0$ for some $r\ne 0$.
	Then $A$ has density at least $\alpha + \frac\delta 2$
	when restricted to some hyperplane.
\end{lemma}

\begin{proof}
	Let $\alpha_0$, \dots, $\alpha_{p-1}$ be the densities of $A$
	in the cosets of the orthogonal subspace $r^\perp$. Then
	\[\delta = \hat{\ind A}(r) = \EE_x\ind A(x)\omega^{-r\cdot x}
	= \frac{\alpha_0 + \alpha_1\omega + \cdots + \alpha_{p-1}\omega^{p-1}}{p}.\]
	By the triangle inequality, we bound
	\begin{align*}
		p\delta &\le
		|\alpha_0 + \alpha_1\omega + \cdots + \alpha_{p-1}\omega^{p-1}|\\
		&\le |(\alpha_0-\alpha) + (\alpha_1-\alpha)\omega+\cdots+
		(\alpha_{p-1}-\alpha)\omega^{p-1}|\\
		&\le |\alpha_0-\alpha| + \cdots + |\alpha_{p-1}-\alpha|\\
		&= \sum_{j=0}^{p-1} |\alpha_j-\alpha| + (\alpha_j-\alpha),
	\end{align*}
	using the fact that $1+\omega+\cdots+\omega^{p-1}=0$
	and applying the triangle inequality.
	Then there exists $j$ such that
	\[|\alpha_j-\alpha| + (\alpha_j-\alpha) \ge \delta,\]
	which gives us the desired bound, since the LHS can only be positive
	if $\alpha_j-\alpha>0$ and so $2\alpha_j-2\alpha\ge\delta$.
\end{proof}

In summary, if $A\subseteq \FF_p^n$ is 3-AP-free and has
$p^n\ge 2\alpha^{-2}$, then we can restrict $A$ to a hyperplane
with respect to which it has density at least $\alpha + \frac{\alpha^2}4$.
Translate everything now so that this hyperplane passes through the origin,
so that we have essentially just reduced the dimension $n$ by one.
\begin{remark}
	Note that this translation relies on the fact that
	our pattern, 3-AP, is translation invariant so that we can safely
	shift the hyperplane to the origin;
	indeed, our proof would break down at this point
	if we were instead trying to forbid e.g. triples satisfying $x+y=z$.
\end{remark}

Suppose now that we repeatedly apply this argument, say $i$ times,
until we have restricted $A$ to a subspace $V_i$ with codimension $i$,
i.e. a copy of $\FF_p^{n-i}$. Let $\alpha_i$ be the density of $A$ in $V_i$,
namely $\alpha_i = |A\cap V_i|/|V_i|$.
From the work above we have the recursive bound
\[\alpha_{i+1}\ge \alpha_i+\frac{\alpha_i^2}{4}.\]
It turns out that going with the crudest possible estimate of
``$\alpha_i$ increases by $\alpha^2/4$ each step'' isn't quite good enough,
so we will be a little less wasteful. Consider instead the time that it takes
for $\alpha_i$ to double. Initially, $\alpha_i$
is increasing by at least $\alpha^2/4$, so it at first takes up to
$\ceil{4/\alpha}$ steps to double. Then, at this point, we can be sure
that $\alpha_i$ increases by at least $\alpha$ per step, so
we need $\ceil{1/\alpha}$ steps to double. In this way we perform at most
\[\sum_{j\le\log_2(1/\alpha)} \ceil{\frac{4^{2-j}}{\alpha}} = O(1/\alpha)\]
steps before the algorithm is forced to terminate
(and even this count is a little lossy in view of the fact that
we can only density increment so long as $2\alpha_i^{-2}\le p^{n-i}$).
Nevertheless, suppose that the algorithm ends after $m$ steps,
leaving us with density $\alpha_m$. Then
\[p^{n-m}\le \alpha_m^{-2}\le\alpha^2,\]
which implies that
\[n\le m+O(\log(1/\alpha))\le O(1/\alpha),\]
where the latter inequality comes from the fact that $m$ is already
at most $O(1/\alpha)$ due to the estimate above;
thus $\alpha = O(1/n)$, or $|A| = O(p^n/n)$ if $A$ is 3-AP-free.
\pagebreak
\section{Act 2: Integers}
\subsection{Fourier stuff speedrun}
In the integers, Fourier analysis is very similar to the things above.
The difference on a visual level is that the transforms
now swap between having domain $\ZZ$ and $\RR/\ZZ$, but
this is only visual: the meat and potatoes of what's happening
is that $\ZZ$ and $\RR/\ZZ$ are duals of one another.
\begin{definition}[Fourier transform, $\ZZ$]
	Given finitely supported $f\colon\ZZ\to\CC$, define
	the Fourier transform $\hat f\colon\RR/\ZZ\to\CC$ by
	\[\hat f(\theta) = \sum_{x\in\ZZ} f(x)e(-x\theta)\]
	where $e(t) := \exp(2\pi it)$, $t\in\RR$.
\end{definition}
This time, we sum in physical space and integrate in frequency space.

\begin{theorem}[Fourier inversion]
	Given finitely supported $f\colon\ZZ\to\CC$, for any $x\in\ZZ$,
	\[f(x) = \int_0^1 \hat f(\theta)e(x\theta)\dif\theta.\]
\end{theorem}

\begin{theorem}[Parseval/Plancherel]
	Given finitely supported $f,g\colon\ZZ\to\CC$,
	\[\sum_{x\in\ZZ}\ol{f(x)}g(x) = \int_0^1\ol{\hat f(x)}\hat g(x)\dif x,\]
	and when $f=g$,
	\[\sum_{x\in\ZZ}|f(x)|^2 = \int_0^1 |\hat f(x)|^2\dif x.\]
\end{theorem}

Convolution again:
\begin{definition}[Convolution]
	Given finitely supported $f,g\colon\ZZ\to\CC$, their convolution $f*g$ is
	\[(f*g)(x) = \sum_{a+b=x} f(a)g(b).\]
\end{definition}

\begin{theorem}[Convolutions become products]
	Given finitely supported $f,g\colon\ZZ\to\CC$ and $\theta\in\RR/\ZZ$,
	\[\hat{f*g}(\theta) = \hat f(\theta)\hat g(\theta).\]
\end{theorem}

\begin{exercise}
	Prove all of the theorems above.
	(Hint: what's the orthonormal basis?)
\end{exercise}

\subsection{Roth's theorem}

Again, the story revolves around the quantity
\[\Lambda(f,g,h) := \sum_{x,y\in\ZZ} f(x)g(x+y)h(x+2y).\]
(Once again we define $\Lambda_3(f) = \Lambda(f,f,f)$.)
\begin{prop}[The $\Lambda$ identity again]
	Given finitely supported $f,g,h\colon\ZZ\to\CC$,
	\[\Lambda(f,g,h) =
	\int_0^1 \hat f(\theta)\hat g(-2\theta)\hat h(\theta)\dif\theta.\]
\end{prop}

\begin{exercise}
	Prove this proposition.
\end{exercise}

With some effort, we can adapt the proof of Roth's theorem in $\FF_p^n$
to a proof in $\ZZ$ --- this is Roth's original proof of his eponymous theorem.

\begin{theorem}[Roth's theorem, mk. II]
	Every 3-AP-free subset of $[N] = \{1,2,\dots,N\}$ has size
	$O(N/\log\log N)$.
\end{theorem}

The main difference in our argument this time is that
since we don't have hyperplanes to restrict to, we will choose to restrict to
subprogressions of a particularly resonant arithmetic progression instead.

Write
\[\norm{\hat f}_\infty := \sup_{\theta}|\hat f(\theta)|,\quad
\norm f_{\ell^2} := \left(\sum_{x\in\ZZ} |f(x)|^2\right)^{1/2}.\]
Then, we have the following version of the counting lemma in $\FF_p^n$:
\begin{lemma}[Counting lemma]
	Let $f,g\colon\ZZ\to\CC$ be finitely supported. Then,
	\[|\Lambda_3(f) - \Lambda_3(g)| \le 3\norm{\hat{f-g}}_\infty
	\max\{\norm f_{\ell^2}^2, \norm g_{\ell^2}^2\}.\]
\end{lemma}

Since we don't have a good analogue $\EE f$ in this case,
we use this version of the counting lemma, which says roughly that
if $f$ and $g$ are close to each other (in the sense of the
Fourier transform of their difference not being too large)
then they have similar 3-AP density.

\begin{proof}
	The idea is to write
	\[\Lambda_3(f) - \Lambda_3(g) =
	\Lambda(f-g,f,f) + \Lambda(g,f-g,f) + \Lambda(g,g,f-g)\]
	and then bound each term in magnitude. The first addend can be bounded
	\begin{align*}
		|\Lambda(f-g,f,f)| &=
		\abs{\int_0^1 \hat{(f-g)}(\theta)\hat f(-2\theta)\hat f(\theta)\dif\theta}\\
		&\le \norm{\hat{f-g}}_\infty
		\abs{\int_0^1\hat f(-2\theta)\hat f(\theta)\dif\theta}\\
		&\le \norm{\hat{f-g}}_\infty
		\left(\int_0^1|\hat f(-2\theta)|^2\dif\theta\right)^{1/2}
		\left(\int_0^1|\hat f(\theta)|^2\dif\theta\right)^{1/2}\\
		&= \norm{\hat{f-g}}_\infty\norm f_{\ell^2}^2.
	\end{align*}
	Similarly, the other two addends can be bounded by
	$\norm{\hat{f-g}}_\infty\norm f_{\ell^2}\norm g_{\ell^2}$ and
	$\norm{\hat{f-g}}_\infty\norm g_{\ell^2}^2$ respectively,
	at which point summing gives the result.
\end{proof}

The other big reason that we want this iteration of the counting lemma
is that it's more fruitful to try estimating the Fourier coefficients
of $\ind A - \alpha\ind{[N]}$, since this time ``excluding zero''
actually looks more like trying to exclude some neighborhood of zero
because our resonant frequencies can be any real number;
it's easier to just shift the function down to have mean zero.

\begin{lemma}[Large Fourier coefficient]
	Let $A\subseteq [N]$ be a 3-AP-free set with $|A| = \alpha N$.
	If $N\ge 5\alpha^{-2}$, then there exists $\theta\in\RR/\ZZ$ such that
	\[|(\ind A - \alpha\ind{[N]})^\wedge(x)| \ge \frac{\alpha^2}{10}N.\]
\end{lemma}

\begin{proof}
	As per the counting lemma we have the estimate
	\[|\Lambda_3(\ind A) - \Lambda_3(\ind{[N]})|\le
	3\norm{(\ind A-\alpha\ind{[N]})^\wedge}_\infty
	\cdot \max\{\norm{\ind A}_{\ell^2},\norm{\alpha\ind{[N]}}_{\ell^2}\}.\]
	We know that $\Lambda_3(\ind A) = |A| =\alpha N$ because $A$ is 3-AP-free.
	We can compute $\Lambda_3(\ind{[N]}) = \floor{N^2/2}+\ceil{N^2/2}\ge N^2/2$
	by counting pairs of integers with same parity (such pairs
	are in bijection with 3-APs, for the middle term of the 3-AP
	is then fixed as the average of the chosen integers).
	Furthermore, the $\ell^2$ norms are respectively
	\[\norm{\ind A}_{\ell^2}^2 =\alpha N,\quad
	\norm{\ind{[N]}}_{\ell^2}^2 = \alpha^2 N,\]
	of which $\alpha N$ is larger, so
	\[\frac{\alpha^3N^2}{2} - \alpha N\le
	\alpha^3\Lambda_3(\ind{[N]}) - \Lambda_3(\ind A) \le
	3\alpha N\norm{(\ind A - \alpha\ind{[N]})^\wedge}_{\infty}.\]
	Therefore,
	\[\norm{(\ind A - \alpha\ind{[N]})^\wedge}_{\infty}
	\ge \frac{1}{6}\alpha^2N - \frac13 \ge \frac{1}{10}\alpha^2 N,\]
	using the assumption that $N\ge 5\alpha^{-2}$,
	which implies the existence of the desired $\theta$.
\end{proof}

Now we want to restrict $A$ to a smaller subspace-like set
with respect to which $A$ has a larger density so that we can
achieve the density increment step. Specifically, we look for
arithmetic progressions whose common difference $d$ makes $d\theta$
reasonably close to an integer.

\begin{lemma}[Dirichlet's lemma]
	Given $\theta\in\RR$ and $0 < \delta < 1$, there exists
	a positive integer $d \le 1/\delta$ such that
	$\norm{d\theta}_{\RR/\ZZ}\le\delta$.
\end{lemma}

\begin{exercise}
	Prove this lemma (it's an easy application of the pigeonhole principle).
\end{exercise}

This lemma tells us about the amount of structure we can impose
in a partition of $[N]$ into arithmetic progressions.
\begin{lemma}
	Let $0 < \eta < 1$ and $\theta\in\RR$. Suppose $N\ge (4\pi/\eta)^6$.
	Then one can partition $[N]$ into arithmetic progressions
	$P_i$ with length between $N^{1/3}$ and $2N^{1/3}$ such that
	\[\sup_{x,y\in P_i} |e(x\theta) - e(y\theta)| < \eta\]
	for each $i$.
\end{lemma}

\begin{proof}[Partition into progressions]
	There exists a positive integer $d < \sqrt N$ such that
	$\norm{d\theta}_{\RR/\ZZ}\le1/\sqrt N$, by Dirichlet's lemma.
	We can then partition $N$ into progressions $P_i$
	with the aforementioned length restriction and common difference $d$
	(exercise: describe how to do this more precisely).
	Then, for any $x$, $y$ in the same progression $P_i$,
	if $x-y = dk$,
	\begin{align*}
		|e(x\theta) - e(y\theta)|\\
		&= |e(dk\theta) - 1|\cdot |e(y\theta)|\\
		&= |e(d\theta) - 1|\cdot|e(d(k-1)\theta) + e(d(k-2)\theta) +\cdots+ 1|\\
		&\le |P_i|\cdot |e(d\theta)-1|\\
		&\le 2N^{1/3}\cdot 2\pi\cdot N^{-1/2}\le \eta,
	\end{align*}
	with $|e(d\theta)-1|$ bounded by the length
	$2\pi\norm{d\theta}_{\RR/\ZZ}$ of the subtended arc on the unit circle.
\end{proof}

Putting all of this together, we have the following density increment:
\begin{lemma}[Density increment]
	Let $A\subseteq [N]$ be 3-AP-free with $|A| = \alpha N$ and
	$N\ge (16/\alpha)^{12}$. Then, there exists a subprogressions
	$P\subseteq [N]$ with $|P|\ge N^{1/3}$ and
	$|A\cap P|\ge (\alpha+\alpha^2/40)|P|$.
\end{lemma}

\begin{proof}
	Let $\theta$ be the frequency corresponding the large Fourier coefficient,
	as found in Lemma 2.11.
	The lower bound on $N$ means that we can apply the partitioning lemma
	with $\eta = \alpha^2/20$, so we split $N$ into $k$ partitions
	$P_1$, \dots, $P_k$ such that for any $x,y\in P_i$ we have
	$|e(x\theta)-e(y\theta)|\le \alpha^2/20$.
	Thus, for each $P_i$,
	\[\abs{\sum_{x\in P_i} (\ind A-\alpha)(x)e(x\theta)}
	\le \abs{\sum_{x\in P_i}(\ind A-\alpha)(x)} + \frac{\alpha^2}{20}|P_i|\]
	(by writing $e(x\theta) = e(m\theta) + (e(x\theta) - e(m\theta))$
	for some fixed $m\in P_i$). Summing over all $i$,
	\begin{align*}
		\frac{\alpha^2}{10}N
		&\le \abs{\sum_{x=1}^N (\ind A-\alpha)(x)e(x\theta)}
		\le \sum_{i=1}^k \abs{\sum_{x\in P_i}(\ind A-\alpha)(x)e(x\theta)}\\
		&\le \sum_{i=1}^k \left(\abs{\sum_{x\in P_i}(\ind A-\alpha)(x)}
		+ \frac{\alpha^2}{20}|P_i|\right)\\
		&= \frac{\alpha^2}{20}N +
		\sum_{i=1}^k \abs{\sum_{x\in P_i}(\ind A-\alpha)(x)},
	\end{align*}
	so in particular
	\[\frac{\alpha^2}{20}\sum_{i=1}^k |P_i|
	\le \sum_{i=1}^k \abs{|A\cap P_i|-\alpha|P_i|}.\]
	Then, we use the same trick of adding
	$|A\cap P_i| -\alpha|P_i|$ to each addend to conclude
	that one of the $|A\cap P_i|-\alpha|P_i|$
	must be at least $\frac{\alpha^2}{40}|P_i|$, as desired.
\end{proof}

We can then translate and scale this progression $P$ so that $A\cap P$
transforms to a subset $A'\subseteq [N']$ with $N' = |P|$,
setting us up to iterate the density increment.

Therefore our setup is as follows:
given $\alpha_0=\alpha$ and $N_0=N$, we can, after $i$ iterations,
reduce to a progression with length $N_i$ in which $A$ has density $\alpha_i$.
As long as $N_i$ is large enough (at least $(16/\alpha_i)^{12}$),
we can apply the increment, giving the recursive relations
\[N_{i+1}\ge N_i^{1/3},\quad \alpha_{i+1}\ge\alpha_i + \frac{\alpha_i^2}{40}.\]
Then, performing a similar doubling analysis, we see that
we can do $m = O(1/\alpha)$ density increments before the algorithm halts.
At this point, since $N_m$ is too small to do any more density increments,
\[N^{1/3^m}\le N_m \le (16/\alpha_m)^{12}\le (16/\alpha)^{12},\]
which implies that
\[N\le (16/\alpha)^{e^{O(1/\alpha)}}\implies \log(\log_{16/\alpha}N) =
\log\log N \le O(1/\alpha) + \log\log(16/\alpha) = O(1/\alpha).\]
Therefore, $\alpha$ is $O(1\log\log N)$,
and so our proof is done!

\begin{remark}
	There is a notable difference in the bound for the $\FF_p^n$ case and here:
	the denominator has an extra layer of log in it;
	this is because the progressions decrease in size
	way faster than the subspaces in the finite field vector space realm.
\end{remark}

\begin{exercise}
	Show that Roth's theorem in $\ZZ$ implies
	that the maximal size of a 3-AP-free set in $\ZZ/N\ZZ$ for odd $N$ is
	$O(N/\log\log N)$.
\end{exercise}

\section{Act 3: Abelian groups with odd order}

\subsection{Fourier once more}
We can also consider arithmetic progressions in arbitrary abelian groups.
Here we will restrict ourselves to just finite abelian groups with odd order.
In this case, frequency space corresponds to the \vocab{Pontryagin dual}
$\hat G$ of $G$, consisting of all homomorphisms $\chi\colon G\to\CC$;
we call these homomorphisms \vocab{characters},
and the characters end up forming the orthonormal basis in this setting.

In this setting we can again define the Fourier transform:
\begin{definition}[Fourier transform, abelian groups]
	Let $f\colon G\to\CC$ be a function. Then its \vocab{Fourier transform}
	$\hat f\colon\hat G\to\CC$ is
	\[\hat f(\chi) := \sum_{x\in G} f(x)\chi(-x).\]
\end{definition}

Then we have the same basic theorems:
\begin{theorem}[Fourier inversion]
	Given a function $\hat f\colon\hat G\to\CC$,
	\[f(x) = \frac{1}{|G|}\sum_{\chi\in\hat G}\hat f(\chi)\chi(x).\]
\end{theorem}

\begin{theorem}[Parseval/Plancherel]
	Given functions $f,g\colon G\to\CC$,
	\[\sum_{x\in G}\ol{f(x)}g(x)
	= \frac1{|G|}\sum_{\chi\in\hat G}\ol{\hat f(\chi)}\hat g(\chi).\]
	In particular, when $f=g$,
	\[\sum_{x\in G}|f(x)|^2 =
	\frac{1}{|G|}\sum_{\chi\in\hat G}|\hat f(\chi)|^2.\]
\end{theorem}

\begin{theorem}[Convolutions become products]
	Given functions $f,g\colon G\to\CC$,
	\[\hat{f*g}(x) = \hat f(x)\hat g(x)\]
	for all $x\in G$.
\end{theorem}

\pagebreak
\subsection{Roth's theorem}

The nice thing about finite abelian groups groups is that we can write
such a group $G$ as a direct sum of finite cyclic groups
\[G \cong \ZZ/k_1\ZZ \oplus \ZZ/k_2\ZZ\oplus\cdots\oplus \ZZ/k_n\ZZ\]
where $k_1\mid \cdots\mid k_n$, so all we really need to do
is to figure out a way to patch together a bunch of cyclic groups
and then borrow a Roth-style bound from the case of $\ZZ$.

Some quick notation:
for an arbitrary finite abelian group $G$,
let $D(G)$ be the maximal cardinality of a 3-AP-free subset $A\subset G$,
and let $c(G)$ be the number $n$ above, i.e. the minimal number
of factors in a decomposition of $G$ into cyclic groups.

\begin{theorem}[Roth's theorem, mk. III]
	For all finite abelian groups $G$ with odd order,
	$D(g)\le 2|G|/c(G)$.
\end{theorem}

Let $d(n)$ be the supremum of $D(G)/|G|$ as $G$ varies across all $G$
such that $c(G)\ge n$. Then, we wish to show that $d(n)\le 2/n$ for all $n$;
we shall do this via induction.
Since $n=1$ and $n=2$ are trivial we shall assume that $c(G)\ge n > 2$.

Let $B = -2A$, and define the functions $f(\chi) = \hat{\ind A}(\chi)$,
$g(\chi) = \hat{\ind B}(\chi)$, and $h(\chi) = d(n-1)|G|\delta(\chi)$,
where $\delta\colon\hat G\to\CC$ is zero everywhere except for at
the trivial character $\chi_0\equiv 1$. Then, we know that
\[\ind A*\ind A*\ind B(0) = |A|\]
is the 3-AP count, which by Fourier inversion gives us the equation
\[|A| = \ind A*\ind A*\ind B(0)
= \frac{1}{|G|}\sum_{\chi\in\hat G}(\ind A*\ind A*\ind B)^\wedge(\chi)
= \frac{1}{|G|}\sum_{\chi\in\hat G}f(\chi)^2g(\chi).\]
On the other hand, noting that $\hat f(\chi_0)$ is the sum of $f$
over all of $G$, we have
\[|A|^2d(n-1) = \frac{1}{|G|}f(\chi_0)^2h(\chi_0)
= \frac{1}{|G|}\sum_{\chi\in\hat G}f(\chi)^2h(\chi).\]
The point here is that $d(n-1)|A|$ should be slightly larger than,
but still quite close to, $1$, and moreover, the difference
between these functions should be pretty close to zero.
This means that we want to prove that $d(n-1)$ and $\ind B(x)$
are ``Fourier-close'' in that their difference $u(x) = d(n-1) - \ind B(x)$
has a Fourier transform $\hat u(\chi) = h(\chi) - g(\chi)$ that is quite small.

We start by getting an upper bound on $\hat u$.

\begin{prop}
	We have
	\[\max_{\chi\in\hat G}|\hat u(\chi)| = d(n-1)|G| - |A|.\]
\end{prop}

\begin{proof}
	We will show that this maximum is achieved if and only if
	$\chi = \chi_0$. Given $\chi\in\hat G$, let $W = \ker\chi$.
	We have that
	\[\ind W*u(x) = \sum_{w\in W}u(x-w) = |W|d(n-1) - |B\cap(x-W)|.\]
	Noting that $|B\cap(x-W)| = |W\cap(x-B)|\le |W|d(n-1)$
	because $W\cap(x-B)$ is 3-AP-free and $c(W)\ge n-1$,
	we can be assured that the expression above is nonnegative.
	
	Then, applying the Fourier transform, we have
	\begin{align*}
		|W||\hat u(\chi)| &= |\hat{\ind W}(\chi)||\hat u(\chi)|
		= |\hat{\ind W*u}(\chi)|\\
		&=\abs{\sum_{x\in G}\ind W*u(x)\chi(-x)}
		\le \sum_{x\in G} \ind W*u(x)\\
		&= |W||G|d(n-1) - \sum_{x\in G} |B\cap (x-W)|\\
		&= |W||G|d(n-1) - |B||W|.
	\end{align*}
	Noting that $|B| = |A|$, this rearranges directly to
	the desired upper bound on the maximum,
	and the upper bound is sharp if and only if all the
	complex numbers $\chi(-x)$ as $x$ varies across $G$
	are scalar multiples of one another, i.e. $\chi = \chi_0$.
\end{proof}

\begin{remark}
	The fact that this bound is sharp at $\chi_0$
	means that $u$ is not particularly resonant with any nontrivial character,
	which is sort of what ``Fourier-close'' should be interpreted as.
\end{remark}

Proposition 3.6 gives us the mechanism for bounding $|A|$:
\begin{align*}
	|G|\cdot |A|(|A|d(n-1) - 1)
	&= \abs{\sum_{\chi\in\hat G}f(\chi)^2(h(\chi) - g(\chi))}\\
	&\le \max_{\chi\in\hat G}|\hat u(x)|\cdot\sum_{\chi\in\hat G}|f(\chi)|^2\\
	&= (|G|d(n-1)-|A|)\cdot |G|\cdot\norm{\ind A}_2^2\\
	&= |G|\cdot |A|(|G|d(n-1) - |A|).
\end{align*}

Since $|G|\inv < 1$ by assumption,
\begin{align*}
	\frac{|A|}{|G|}&\le \frac{d(n-1) + |G|\inv}{d(n-1)+1}
	\le \frac{d(n-1) + 3^{-n}}{d(n-1) + 1}\\
	&\le \frac{2/(n-1) + 3^{-n}}{2/(n-1) + 1} < \frac 2n
\end{align*}
where the second-to-last inequality comes from
\[x\mapsto \frac{x+a}{x+b}\]
being increasing on $(0,+\infty)$ if $0 \le a < b$
and the last one from some simple algebraic manipulation and induction.
This completes the proof of this iteration of Roth's theorem!\\

This theorem can also give us a more concrete bound on $|A|$ in terms of $|G|$.
For example, if we take for granted that $D(\ZZ/m\ZZ) = O(m/(\log m)^\alpha)$
for some constant $\alpha$ (as was proven by Szemer\'edi and Heath-Brown
in the late 80s), then we show that $D(G) = O(|G|/(\log |G|)^\beta)$
for $\beta = \alpha/(1+\alpha)$. If $c(G)$ is large enough,
namely $c(G) > (\log |G|)^\beta$, then the work above finishes,
so suppose otherwise with $t = c(G) \le (\log |G|)^\beta$.

Letting $A\subseteq G$ be 3-AP-free, let $H$ be a cyclic subgroup of $G$
such that $|H|\ge |G|^{1/t}$ (say, the factor in the decomposition of $G$
with maximal order). There exists a coset $x+H$ such that
the density of $A$ with respect to this coset is at least that of $A$ in $G$:
\[|A\cap(x+H)|\ge \frac{|A||H|}{|G|}.\]
However, the $A\cap(x+H)$ is 3-AP-free in a cyclic group,
so it is susceptible to the standard Roth's theorem,
where we can say that
\[|A\cap (x+H)| = O(|H|/(\log |H|)^\alpha).\]
Therefore
\[|A| = O\left(\frac{|G|}{(\log |H|)^\alpha}\right)
= O\left(\frac{|G|t^\alpha}{(\log |G|)^\alpha}\right)
= O\left(\frac{|G|}{(\log |G|)^\beta}\right),\]
as desired.




















\end{document}